\section{evnx vault}

\begin{frame}{\cmd{evnx vault} — the unit of sharing}
  \begin{block}{Synopsis}\texttt{evnx vault <SUBCOMMAND>}\end{block}

  \begin{tabular}{@{}ll@{}}
    \toprule
    \texttt{vault create <NAME>} & \texttt{app} or \texttt{app/production} \\
    \texttt{vault list}          & Vaults you can reach \\
    \texttt{vault delete <NAME>} & Soft delete \\
    \texttt{vault share <NAME> --with <EMAIL>} & Wrap the key for someone else \\
    \texttt{vault members <NAME>} & Who has access, at what role \\
    \texttt{vault role <NAME> --user --role} & Change a role \\
    \texttt{vault revoke <NAME> --user <EMAIL>} & Remove \textbf{and re-key} \\
    \bottomrule
  \end{tabular}

  \vspace{0.4em}
  \begin{alertblock}{Two names for the same thing}
    \cmd{vault share} takes \flag{--with}; \cmd{vault revoke} and
    \cmd{vault role} take \flag{--user}. Same concept, two flags --- worth
    knowing before you guess.
  \end{alertblock}

  \vspace{0.3em}
  {\footnotesize A vault name is \texttt{name} or \texttt{name/environment};
  \texttt{app/production} binds naturally to \file{.env.production}.}
\end{frame}

\begin{frame}{Two wrap modes — and the distinction is load-bearing}
  \begin{tabular}{@{}p{0.26\textwidth}p{0.40\textwidth}p{0.24\textwidth}@{}}
    \toprule
    \textbf{Whose copy} & \textbf{How it is wrapped} & \textbf{Post-quantum} \\
    \midrule
    The creator's own & XChaCha20 under an HKDF subkey of the master key & \evnxCheck{} symmetric only \\
    A member's, shared & \textbf{hybrid X25519 + ML-KEM-768} & \evnxCheck{} holds if either survives \\
    \bottomrule
  \end{tabular}

  \vspace{0.8em}
  The creator's copy deliberately does \emph{not} go through ECDH — that would
  put every solo vault behind X25519 and expose it to
  harvest-now-decrypt-later.

  \vspace{0.4em}
  \begin{block}{There is no classical-only path}
    Not in the crate, not in the CLI, and not in the database, where a CHECK
    constraint makes half a wrap unstorable. Sharing with an account that has no
    ML-KEM key on file is \textbf{refused}, not downgraded.
  \end{block}
\end{frame}

\begin{frame}[fragile]{Revocation actually revokes}
  \capturepart{63-vault-revoke-help.txt}{1}{8}

  {\footnotesize This is the best-written help in the CLI because it states the
  \emph{limit} of what the command can do — the sentence most security tooling
  omits.}
\end{frame}

\begin{frame}{The honest caveat}
  \begin{alertblock}{The recipient's public keys come from the server}
    A malicious server could substitute its own and read what you share. There
    is no third party to check against, and out-of-band fingerprint verification
    is not built.
  \end{alertblock}

  \vspace{0.6em}
  So \emph{``the server cannot read your secrets''} becomes \emph{``\ldots cannot
  read them \textbf{passively}''} the moment you share.

  \vspace{0.8em}
  \textbf{Use cases}
  \begin{itemize}\footnotesize
    \item \texttt{evnx vault create app/production}
    \item \texttt{evnx vault share app/production --with teammate@example.com}
    \item \texttt{evnx vault role app/production --user t@x.com --role viewer}
    \item \texttt{evnx vault revoke app/production --user former@example.com}
  \end{itemize}
\end{frame}
